% Generated draft: requires mathematical and source review.
\documentclass[11pt]{article}
\usepackage[margin=25mm]{geometry}
\usepackage{fontspec}
\setmainfont{Latin Modern Roman}
\usepackage{amsmath,amssymb,mathtools,unicode-math}
\setmathfont{Latin Modern Math}
\usepackage{enumitem,needspace,xurl,hyperref}
\setlength{\parindent}{0pt}
\setlength{\parskip}{5pt}
\setlength{\emergencystretch}{3em}
\newcommand{\N}{\mathbb{N}}
\newcommand{\Z}{\mathbb{Z}}
\newcommand{\Q}{\mathbb{Q}}
\newcommand{\R}{\mathbb{R}}
\newcommand{\C}{\mathbb{C}}
\newcommand{\F}{\mathbb{F}}
\newcommand{\A}{\mathbb{A}}
\newcommand{\PP}{\mathbb{P}}
\newcommand{\E}{\mathbb{E}}
\newcommand{\eps}{\varepsilon}
\newcommand{\dd}{\,\mathrm{d}}
\begin{document}
\section*{UnsolvedMath ID: 30001227}
\textbf{Source problem code:} catalogue entry
\subsection{Definitions and mathematical statement}
Let \(n\geq 3\) be an integer and let \(H\) be a Hilbert space. Write \(\mathcal B(H)\) for the bounded linear operators on \(H\). A commuting tuple of contractions is a tuple \(T=(T_1,\ldots,T_n)\in\mathcal B(H)^n\) satisfying
\[
\|T_i\|\leq 1\quad(1\leq i\leq n),\qquad T_iT_j=T_jT_i\quad(1\leq i,j\leq n).
\]
The supplied problem asks for an effective model for arbitrary such tuples, analogous to Ando's model for two commuting contractions.

The supplied record does not define an ``effective model,'' identify the particular formulation of Ando's model intended, or specify which features of that model should extend to \(n\geq3\). Thus this is a model-construction problem whose required output and success criteria remain unspecified; it does not yet determine a precise theorem to prove. In particular, no specific dilation, representation, or equivalence requirement is imposed by the supplied statement. The scalar field of \(H\) and any additional hypotheses on the operators are also unspecified. The finite-tuple formulation above expresses the request for each integer \(n\geq3\); the record does not address infinite families.
\subsection{Short English statement}
Find an effective model for three or more commuting contractions on a Hilbert space, analogous to Ando's model for two commuting contractions. The intended meaning of the model and the features it must preserve require clarification. The supplied record labels the problem open; this is not an independently verified status.
\subsection{Sources}
\begin{itemize}
\item Supplied local registry record: permanent UnsolvedMath ID \(30001227\), source problem code \texttt{catalogue entry}, title ``Models for Three or More Commuting Contractions.'' The record gives proposed year 2009 and does not identify an individual proposer.
\item Supplied source citation: \emph{Hilbert Modules and Complex Geometry} (2009), Oberwolfach Reports, pp.\ 1047--1100, DOI 10.4171/owr/2009/20. Source URL and supplied literature source labeled ``original report'': \url{https://doi.org/10.4171/owr/2009/20}. The supplied background describes Oberwolfach Reports as documenting mathematical workshops and their open-problem discussions.
\item Supplied literature assessment, checked \texttt{2026-08-21}: status ``open,'' with the assessment that the report explicitly leaves the effective model problem for three or more commuting contractions open. This is supplied dated curation derived from the Oberwolfach Reports record, not an independently verified literature assessment; time-sensitive status and cited claims require independent verification.
\item The registry marks the statement as ``exact'' and supplies a manual-review assertion that the extracted statement is self-contained, mathematically coherent, and free of LaTeX or report-layout artifacts. That supplied assertion does not resolve the unspecified model requirements and is not independent mathematical review of this draft.
\end{itemize}
\end{document}
